
Declaramos um vetor de população com m=20 cadeias de 16 bits randomizadas.A aptidão foi escolhida como $f(x) = 1/(w_{13}+600)$, que é positiva dentro de todo o espaço de busca e máxima no ponto de mínimo global da função. O espaço de busca foi limitado a $-500 \leq x \leq 500$ e $-500 \leq y \leq 500$. As constantes pertinentes a mutação e crossover foram escolhidas iguais a 0.01 e 0.75 respectivamente.\\

Rodando o algoritmo com um número de iterações igual a 36, que é suficiente para garantir a convergência dos membros da população para dentro da bacia de mínimo, obtemos o valor da posição do mínimo global como (- 132.71815, - 425.27714), e o valor da função $w_{13}$ para este ponto é -529.23716.\\

Podemos ter uma confiança maior que este é um mínimo local da função ao invocar a função $w_{13}$ para valores de argumentos próximo à posição do mínimo encontrado, embora este teste não seja exaustivo e não exista garantia de que o ponto é realmente o mínimo da bacia encontrada. De fato, o algoritmo da nuvem de partículas encontra um valor menor para $w_{13}$, embora deva ser notado que a população declarada para a nuvem de partículas é muito maior e o algoritmo rodou durante mais iterações. Apesar dos resultados convergirem não conseguimos cumprir o objetivo expresso do trabalho de comparar os dois algoritmos de forma justa, pois seria necessário afinar os parâmetros de cada um para que consumissem aproximadamente o mesmo custo computacional. \\ 